\section{Bölüm sonu çözümlü problemler}

\begin{enumerate}
  \item \textbf{Tek örneklem testi.} $n=16$, $\bar x=54$, $s=8$ ve $\mu_0=50$ için $t$ değerini ve serbestlik derecesini bulun.

  \textbf{Çözüm:} $SE=8/\sqrt{16}=2$ ve $t=(54-50)/2=2$'dir. Serbestlik derecesi $16-1=15$ olur. Karar için alternatif hipotez ve $t(15)$ dağılımından $p$-değeri gerekir.

  \item \textbf{Welch testi.} İki grubun ortalamaları 40 ve 46; standart sapmaları 6 ve 8; örneklem hacimleri 25 ve 16'dır. Farkı ve Welch standart hatasını bulun.

  \textbf{Çözüm:} Birinci eksi ikinci farkı $40-46=-6$'dır. $SE=\sqrt{6^2/25+8^2/16}=\sqrt{1{,}44+4}=\sqrt{5{,}44}\approx2{,}33$ ve $t\approx-6/2{,}33=-2{,}57$ olur.

  \item \textbf{Eşleştirilmiş test.} On öğrencinin son-test eksi ön-test fark ortalaması $-3$, fark standart sapması 4'tür. Test istatistiğini ve $d_z$ değerini bulun.

  \textbf{Çözüm:} $SE=4/\sqrt{10}\approx1{,}265$ ve $t=-3/1{,}265\approx-2{,}37$'dir; $sd=9$ olur. $d_z=\bar d/s_D=-3/4=-0{,}75$'tir. Negatif yön, son-test puanlarının daha düşük olduğunu gösterir.
\end{enumerate}